\documentclass[12pt, twocolumn, landscape, a4paper]{article}

\usepackage[margin=1.2cm]{geometry}
\usepackage[utf8]{inputenc}
\usepackage[spanish,es-noshorthands]{babel}
\usepackage{amssymb, amsmath, amsbsy}
\usepackage{verbatim}
\usepackage{mathpazo}
\usepackage{graphicx}

\usepackage{tikz}
\usetikzlibrary{angles,quotes}

\usepackage[pdftex]{hyperref}
\hypersetup{ pdftitle = {Ejercicios en vectores, geometría},pdfauthor = {Luis Carrillo Gutiérrez} }

\renewcommand{\familydefault}{\sfdefault}
\pagestyle{empty}

\begin{document}
\paragraph*{Vectores -- Ejercicios}

\begin{itemize}
\item{Determinar el valor del parámetro ``\scalebox{1.1}{$\,k\,$}'' para que los vectores :
$$
\vec{x} = k\,\vec{u} - 2\,\vec{v} + 3\,\vec{w}
$$
$$
\vec{y} = - \vec{u} + k\,\vec{v} + \vec{w}
$$
sean: ortogonales ($\vec{x}\perp\vec{y}$),\qquad paralelos ($\vec{x}\parallel\vec{y}$).
}
\item{Dados los puntos $A(1,\,0,\,1),\quad B(1,\,1,\,1)$ y\quad $C(1,\,6,\,\mu)$, se pide :
\begin{itemize}
\item{Hallar los valores del parámetro ``\scalebox{1.1}{$\mu$}'' para que los puntos estén alineados.}
\item{Hallar si existen valores de ``\scalebox{1.1}{$\mu$}'' para los cuales $A$, $B$ y $C$ sean los tres vértices de un paralelogramo de área 3 y en caso afirmativo, calcular}
\end{itemize}
}
\item{Hallar un vector perpendicular a $\overrightarrow{\text{U}} = (2,\,3,\,4)$ y $\vec{\text{V}} =(-1,\,3,\,-5)$, y que sea unitario.}
\item{Hallar dos vectores de módulo la unidad y ortogonales a $(2,\,-2,\,3)$ y $(3,\,-3,\,2)$.}
\end{itemize}

% Comparing two vectors by Glenn Research Center.
\paragraph*{Comparar vectores} -- \\

Un vector tiene {\tt magnitud} y {\tt dirección}.
% vector quantity ???
\begin{itemize}
\item{Si un vector $\overrightarrow{A}$ y un vector $\overrightarrow{B}$, tienen la misma dirección pero con diferentes magnitudes, entonces $\overrightarrow{A} \neq \overrightarrow{B}$.

\begin{tikzpicture}
	\begin{scope}[-latex]
	\draw (0, 0) -- (1, 3);
	\draw (1, 0) -- (1.5, 1.4);
	\end{scope}
	\draw (.1, 1.45) node{$\vec{A}$};
	\draw (1.575, .75) node{$\vec{B}$};
\end{tikzpicture}}

\item{Si un vector $\overrightarrow{A}$ y un vector $\overrightarrow{B}$, tienen la misma magnitud pero con diferente dirección, entonces $\overrightarrow{A} \neq \overrightarrow{B}$.

\begin{tikzpicture}
	\begin{scope}[-latex]
	\draw (0, 0) -- (1, 3);
	\draw (1, 0) -- (2.5, 1.4);
	\end{scope}
	\draw (.1, 1.45) node{$\vec{A}$};
	\draw (2.1, .5) node{$\vec{B}$};
\end{tikzpicture}}

\item{Si un vector $\overrightarrow{A}$ y un vector $\overrightarrow{B}$, tienen la misma dirección y la misma magnitud, entonces $\overrightarrow{A} = \overrightarrow{B}$.

\begin{tikzpicture}
	\begin{scope}[-latex]
	\draw (0, 0) -- (1, 3);
	\draw (1, 0) -- (2, 3);
	\end{scope}
	\draw (.1, 1.45) node{$\vec{A}$};
	\draw (1.76, 1.45) node{$\vec{B}$};
\end{tikzpicture}}
\end{itemize}
\end{document}
